
\documentclass{article}
\usepackage[utf8]{inputenc}
\usepackage[portuguese]{babel}
\usepackage{booktabs}
\usepackage{hyperref}

\title{Análise Linguística: https://pt.wikipedia.org/wiki/Literatura\_fant\%C3\%A1stica}
\author{Gonçalo Magalhães}
\date{\today}

\begin{document}

\maketitle

\begin{abstract}
Frases mais relevantes extraídas através do modelo n-grams:
\begin{itemize}
\item De acordo com a autora, na literatura fantástica européia, ao contrário da produzida na América Latina, há uma preocupação em preservar o real quando algo sobrenatural ocorre, mesmo que a explicação apareça apenas no desfecho da obra.
\item Já no terceiro capítulo, a autora passa a conceituar as diversas nomenclaturas utilizadas quando se refere à literatura fantástica, como o realismo mágico, o maravilhoso e o alegórico, que são posteriormente, no capítulo seguinte, utilizadas como apoio para comparar-se o Fantástico produzido na Europa com o Fantástico da “Hispano-América e Brasil ”.
\item A autora, logo no capítulo inicial, apóia-se em textos de dois autores consagrados da literatura fantástica, E. T. A. Hoffmann , Laura Esquível e Gabriel García Márquez , para conceituar e explicitar as semelhanças e as diferenças dos textos que, apesar de ambos pertencerem à literatura fantástica, apresentam características peculiares que os enquadram em diferentes concepções do gênero.
\end{itemize}
\end{abstract}

\section{Entidades Nomeadas (NER)}
\begin{itemize}
\item \textbf{PER}: J.R.R. Tolkien
\item \textbf{MISC}: L   • 
\item \textbf{MISC}: Saga The Witcher
\item \textbf{MISC}: Wikipédia
\item \textbf{PER}: Todorov
\item \textbf{LOC}: Fantasia
\item \textbf{PER}: Luís Filipe Silva
\item \textbf{LOC}: América
\item \textbf{MISC}: Dark Fantasy
\item \textbf{ORG}: Internacional
\item \textbf{LOC}: Fantasia Histórica
\item \textbf{PER}: Clive Barker
\item \textbf{MISC}: ABW
\item \textbf{MISC}: O Senhor dos Anéis
\item \textbf{PER}: João Barreiros
\item \textbf{LOC}: Letras
\item \textbf{MISC}: A Jangada de Pedra
\item \textbf{ORG}: Especulativa
\item \textbf{PER}: Tzvetan Todorov
\item \textbf{MISC}: As Crônicas de Nárnia
\item \textbf{MISC}: Mostra Curta Fantástico
\item \textbf{MISC}: Realismo Mágico
\item \textbf{PER}: Jorge Luis Borges
\item \textbf{MISC}: Fantasia Sombria
\item \textbf{PER}: Selma Rodrigues
\item \textbf{PER}: D. Afonso Henriques
\item \textbf{PER}: Guy Gavriel Kay
\item \textbf{PER}: Arvède Barine
\item \textbf{MISC}: “Introdução à Literatura Fantástica
\item \textbf{PER}: Neil Gaiman
\item \textbf{LOC}: Cinefantasy
\item \textbf{ORG}: América Latina
\item \textbf{MISC}: Científica
\item \textbf{MISC}: The Witcher
\item \textbf{PER}: João de Melo
\item \textbf{MISC}: Iluminismo
\item \textbf{ORG}: Instrumentos Mortais
\item \textbf{PER}: Joel Perpétuo
\item \textbf{LOC}: Ajude
\item \textbf{LOC}: cidade de São Paulo
\item \textbf{MISC}: Ir
\item \textbf{MISC}: N   • 
\item \textbf{PER}: Selma Calasans Rodrigues
\item \textbf{PER}: Gabriel García Márquez
\item \textbf{PER}: Lídia Jorge
\item \textbf{ORG}: Festival Internacional de Cinema Fantástico CINEFANTASY
\item \textbf{MISC}: Literatura Fantástica
\item \textbf{PER}: Laura Esquível
\item \textbf{LOC}: Portugal
\item \textbf{MISC}: A Mostra Curta Fantástico
\item \textbf{MISC}: Rota do Passado
\item \textbf{ORG}: TODOROV
\item \textbf{LOC}: Terramoto de 1755
\item \textbf{PER}: C. S. Lewis
\item \textbf{MISC}: Inscritos
\item \textbf{MISC}: História
\item \textbf{ORG}: Google
\item \textbf{MISC}: Realismo do século XIX
\item \textbf{MISC}: Inglês
\item \textbf{PER}: Vampire Academy
\item \textbf{MISC}: dark fantasy
\item \textbf{LOC}: Ficção Científica
\item \textbf{MISC}: Encontre
\item \textbf{LOC}: Europa
\item \textbf{PER}: Evolução Contemporânea
\item \textbf{LOC}: Selma
\item \textbf{ORG}: Cinefantas
\item \textbf{PER}: Mary Gentle
\item \textbf{PER}: Andrzej Sapkowski
\item \textbf{LOC}: Mostra
\item \textbf{MISC}: Subgéneros
\item \textbf{PER}: José Saramago
\item \textbf{LOC}: Em Portugal
\item \textbf{MISC}: Realismo Mágico: Consolidado por Gabriel García Márquez
\item \textbf{LOC}: Subgêneros
\item \textbf{MISC}: Fantástico da “Hispano-América e Brasil
\item \textbf{PER}: Júlio Verne
\item \textbf{MISC}: Fantástico
\item \textbf{MISC}: WP
\item \textbf{LOC}: Universidade Federal do Rio de Janeiro
\item \textbf{PER}: T. A. Hoffmann
\item \textbf{MISC}: Origem: Wikipédia
\item \textbf{LOC}: Mostras
\item \textbf{PER}: Margaret Atwood
\item \textbf{LOC}: Brasil
\item \textbf{PER}: Isabel Allende
\end{itemize}
\clearpage

\section{Análise de Padrões Gramaticais}
\begin{table}[h]
\centering
\begin{tabular}{|l|p{8cm}|}
\hline
\textbf{Tipo de Padrão} & \textbf{Trecho Extraído} \\ \hline
OBJETO & há uma preocupação \\ \hline
OBJETO & preservar o real \\ \hline
ACAO & sobrenatural ocorre \\ \hline
MOVIMENTO & refere à literatura \\ \hline
MOVIMENTO & utilizadas como apoio \\ \hline
OBJETO & comparar-se o Fantástico \\ \hline
MOVIMENTO & apóia-se em textos \\ \hline
MOVIMENTO & consagrados da literatura \\ \hline
OBJETO & explicitar as semelhanças \\ \hline
MOVIMENTO & pertencerem à literatura \\ \hline
OBJETO & apresentam características \\ \hline
MOVIMENTO & enquadram em diferentes concepções \\ \hline

\end{tabular}
\end{table}

\clearpage

\begin{thebibliography}{9}
\bibitem{fonte} Wikipédia, \textit{https://pt.wikipedia.org/wiki/Literatura\_fant\%C3\%A1stica}, disponível em \url{https://pt.wikipedia.org/wiki/Literatura_fant%C3%A1stica}.
\end{thebibliography}

\end{document}
